%%%%%%%%%%%%%%%%%%%%%%%%%%%%%%%%%%%%%%%%%%%%%%%%%%%%%%%%%%%%%%%%%%%%%%%%
%                                                                      %
%     File: Thesis_Abstract.tex                                        %
%     Tex Master: Thesis.tex                                           %
%                                                                      %
%     Author: Andre C. Marta                                           %
%     Last modified :  2 Jul 2015                                      %
%                                                                      %
%%%%%%%%%%%%%%%%%%%%%%%%%%%%%%%%%%%%%%%%%%%%%%%%%%%%%%%%%%%%%%%%%%%%%%%%

\section*{Abstract}

\addcontentsline{toc}{section}{Abstract}

\noindent The development of control algorithms for autonomous vehicles at extreme environments has attracted the attention of researchers in the last decades, mainly due to the need of improving the wheeled locomotion efficiency, but also to increase the lifetime of these systems. Firstly, this work presents a safe path generation algorithm, that avoids very uneven or highly inclined locations. Likewise, several traversability maps variants are generated by a nonlinear optimization problem, that take into account the vehicle's geometry and adaptability to the terrain, as well as the elevation and terrain roughness. Then, the \ti{fast marching} method is used to yield an unique, global minimum potential field, that contains directions to a given goal in every position, with the aim to avoid high cost of traversal locations. Lastly, the \texttt{acados} library provides known real-time nonlinear model predictive methods which are implemented to safely guide the wheeled robot through irregular terrain. Apart from following a reference, the controller also enforces traction constraints in order to minimize slippage, by taking advantage of the Coulomb friction model theory. Moreover, the vehicle's collision avoidance is ensured by reformulating the initial provided reference, such that it does not cross obstacles, but also by assigning to each node of the control horizon a convex, free of obstacles area.

\vfill

\textbf{\Large Keywords:} Uneven terrain, Traversability map, Optimal path, Real-time control, Nonlinear Model Predictive control, Low slip and collision avoidance.

